\documentclass[11pt, a4paper]{article}

% Gemeinsame Style-Definitionen (TEXINPUTS muss templates/ enthalten — siehe scripts/build_tex.py)
% bfw-style.tex — gemeinsame Style-Definitionen für alle Handouts/Aufgabenhefte
% Voraussetzung: Kompilation mit xelatex (wegen fontspec)

\usepackage[ngerman]{babel}
\usepackage{geometry}
\geometry{a4paper, top=2.2cm, bottom=2.2cm, left=2.3cm, right=2.3cm}

\usepackage{fontspec}
% Fallback-Kette: Source Sans Pro -> Calibri -> Segoe UI -> Latin Modern Sans
% (Latin Modern ist immer mit MiKTeX vorhanden)
\IfFontExistsTF{Source Sans Pro}{%
  \setmainfont{Source Sans Pro}[Ligatures=TeX]%
}{\IfFontExistsTF{Calibri}{%
  \setmainfont{Calibri}[Ligatures=TeX]%
}{\IfFontExistsTF{Segoe UI}{%
  \setmainfont{Segoe UI}[Ligatures=TeX]%
}{%
  \setmainfont{Latin Modern Sans}[Ligatures=TeX]%
}}}
\IfFontExistsTF{Source Code Pro}{%
  \setmonofont{Source Code Pro}[Scale=0.92]%
}{\IfFontExistsTF{Consolas}{%
  \setmonofont{Consolas}[Scale=0.92]%
}{%
  \setmonofont{Latin Modern Mono}[Scale=0.92]%
}}

\usepackage{xcolor}
% Farbpalette: hell, freundlich, modern
\definecolor{bfwPrimary}{HTML}{0EA5A4}    % Petrol/Türkis - Primär
\definecolor{bfwPrimaryDark}{HTML}{0F766E} % dunkler Petrol für Headings
\definecolor{bfwAccent}{HTML}{F59E0B}     % Amber - Akzent
\definecolor{bfwSuccess}{HTML}{10B981}    % Grün - Erfolg / Lösung
\definecolor{bfwWarn}{HTML}{F97316}       % Orange - Achtung
\definecolor{bfwInk}{HTML}{0F172A}        % fast Schwarz
\definecolor{bfwInkSoft}{HTML}{475569}    % gedämpftes Grau
\definecolor{bfwBgSoft}{HTML}{F8FAFC}     % off-white
\definecolor{bfwBgTeal}{HTML}{ECFEFF}     % helles Türkis
\definecolor{bfwBgAmber}{HTML}{FEF3C7}    % helles Gelb
\definecolor{bfwBgRose}{HTML}{FFE4E6}     % helles Rosé für Eselsbrücken

\usepackage[colorlinks=true, linkcolor=bfwPrimaryDark, urlcolor=bfwPrimaryDark]{hyperref}
\usepackage{enumitem}
\setlist[itemize]{leftmargin=*, itemsep=2pt, topsep=2pt}
\setlist[enumerate]{leftmargin=*, itemsep=2pt, topsep=2pt}

\usepackage[most]{tcolorbox}
\usepackage{booktabs}
\usepackage{tikz}
\usetikzlibrary{decorations.pathmorphing, positioning, shapes.geometric, arrows.meta}

% Merkkästchen — wichtige Definition / Kernpunkt
\newtcolorbox{merksatz}[1][]{%
  enhanced, breakable,
  colback=bfwBgTeal,
  colframe=bfwPrimary,
  boxrule=0.6pt,
  arc=4pt,
  left=10pt, right=10pt, top=8pt, bottom=8pt,
  fonttitle=\bfseries\color{bfwPrimaryDark},
  title={\faLightbulb\ Merksatz},
  #1
}

% Eselsbrücke — Mnemoniks
\newtcolorbox{eselsbruecke}[1][]{%
  enhanced, breakable,
  colback=bfwBgRose,
  colframe=bfwAccent,
  boxrule=0.6pt,
  arc=4pt,
  left=10pt, right=10pt, top=8pt, bottom=8pt,
  fonttitle=\bfseries\color{bfwWarn},
  title={Eselsbrücke},
  #1
}

% Beispiel-Box
\newtcolorbox{beispiel}[1][]{%
  enhanced, breakable,
  colback=bfwBgSoft,
  colframe=bfwInkSoft,
  boxrule=0.4pt,
  arc=3pt,
  left=10pt, right=10pt, top=8pt, bottom=8pt,
  fonttitle=\bfseries\color{bfwInk},
  title={Beispiel},
  #1
}

% Lösungs-Box (für Aufgabenhefte)
\newtcolorbox{loesung}[1][]{%
  enhanced, breakable,
  colback=bfwBgAmber!40,
  colframe=bfwSuccess,
  boxrule=0.6pt,
  arc=3pt,
  left=10pt, right=10pt, top=8pt, bottom=8pt,
  fonttitle=\bfseries\color{bfwSuccess!70!black},
  title={Lösung},
  #1
}

% Glossar-Eintrag
\newcommand{\glossareintrag}[2]{%
  \par\noindent\textbf{\color{bfwPrimaryDark}#1} \enspace #2\par\smallskip
}

% Section-Stil — etwas farbiger als Standard
\usepackage{titlesec}
\titleformat{\section}
  {\Large\bfseries\color{bfwPrimaryDark}}
  {\thesection}{0.6em}{}
\titleformat{\subsection}
  {\large\bfseries\color{bfwPrimary}}
  {\thesubsection}{0.5em}{}
\titleformat{\subsubsection}
  {\normalsize\bfseries\color{bfwInk}}
  {\thesubsubsection}{0.4em}{}

% TikZ-Stil "skizzenhaft" — etwas weicher als Rough.js, aber stabil
\tikzset{
  roughbox/.style = {
    draw=bfwPrimaryDark, thick, fill=bfwBgTeal,
    rounded corners=4pt, inner sep=6pt
  },
  roughaccent/.style = {
    draw=bfwAccent, thick, fill=bfwBgAmber,
    rounded corners=4pt, inner sep=6pt
  },
  roughline/.style = {
    draw=bfwInkSoft, thick, -{Stealth[length=2.5mm]}
  }
}

% Bullet-Symbole (bewusst kein FontAwesome — vermeidet Font-Installations-Probleme)
\newcommand{\faLightbulb}{$\bullet$}
\newcommand{\faBookmark}{$\bullet$}

% Header/Footer
\usepackage{fancyhdr}
\pagestyle{fancy}
\fancyhf{}
\renewcommand{\headrulewidth}{0pt}
\fancyfoot[L]{\small\color{bfwInkSoft}\thepage}
\fancyfoot[R]{\small\color{bfwInkSoft} BFW M-064 Beschaffung im Büromanagement}


\title{\vspace*{-1.5cm}\Huge\bfseries\color{bfwPrimaryDark}Tag~11: Postbearbeitung --- Postausgang\\[0.4em]
  \large\color{bfwInkSoft}\textnormal{Von der Abteilung zur Zustellung --- versandfertig in sieben Schritten}}
\author{\textsf{Büroprozesse gestalten und Arbeitsvorgänge organisieren \textperiodcentered{} M-067}\\
\textsf{\small\copyright\ Can Siebert\,\textperiodcentered\,2026}}
\date{13.07.2026}

\begin{document}
\maketitle
\thispagestyle{empty}

\vspace{1em}
\begin{merksatz}[title={\bfseries Worum geht es heute?}]
Jeder Brief, der ein Unternehmen verlässt, durchläuft einen festen Ablauf: kontrollieren,
sortieren, falzen, kuvertieren, wiegen, frankieren und aufgeben. Diese Schritte
entscheiden über Kosten, Tempo und Rechtssicherheit der Ausgangspost. In diesem Handout
lernen Sie den vollständigen Postausgangsprozess kennen --- von den Versand- und
Sendungsarten über Falzen und Frankieren bis zur Poststraße. Außerdem ordnen Sie die
moderne Praxis ein: Online-Frankierung, Versandsoftware, Hybridpost und die elektronische
Post (E-Mail, De-Mail, EGVP, elektronische Signatur). Das ist Grundlagenwissen für die
betriebliche Praxis und die IHK-Abschlussprüfung.
\end{merksatz}

\tableofcontents
\vspace{1em}
\hrule
\vspace{1em}

\section{Der Ablauf des Postausgangs}

Ist in einem Unternehmen eine \emph{Poststelle} vorhanden, so wird die in den einzelnen
Abteilungen gesammelte Ausgangspost zum Versenden dorthin gebracht. In der Poststelle
fallen verschiedene Arbeitsschritte an, um die Post \emph{versandfertig} zu machen. Diese
Schritte laufen in einer festen Reihenfolge ab und bauen aufeinander auf.

\begin{enumerate}
  \item \textbf{Kontrollieren} --- Prüfung auf Unterschrift, Adresse und Vollständigkeit der Anlagen.
  \item \textbf{Sortieren} --- nach Sendungsart, Versandart, Gewicht sowie In- und Auslandssendungen.
  \item \textbf{Falzen} --- Falten der Schriftstücke ins passende, möglichst kostengünstige Format.
  \item \textbf{Kuvertieren und Verschließen} --- Einstecken in geeignete Umschläge und Verschließen.
  \item \textbf{Wiegen} --- Ermitteln des Gewichts als Grundlage für das Porto.
  \item \textbf{Frankieren} --- Freimachen, also Bezahlen des Portos.
  \item \textbf{Aufgeben} --- Übergabe an den Versanddienstleister.
\end{enumerate}

\begin{eselsbruecke}
\textbf{„Karl Sortiert Flink, Kuvertiert, Wiegt, Frankiert, Aufgibt.”} --- die Anfangs\-buch\-staben
\emph{Ko-So-Fa-Ku-Wi-Fr-Auf} erinnern an die sieben Arbeitsschritte des Postausgangs in
der richtigen Reihenfolge.
\end{eselsbruecke}

\subsection{Kontrollieren und Sortieren}

Zuerst wird die Ausgangspost \emph{kontrolliert}: Ist der Brief unterschrieben? Ist die
Adresse vorhanden? Sind alle Anlagen vollständig? Wird die Sendung nicht in einem
\emph{Fensterbriefumschlag} verschickt und fehlt die Adresse, so muss sie durch eine
Adressiermaschine oder einen Adressaufkleber ergänzt werden.

Anschließend wird die Post \emph{sortiert} --- nach \emph{Sendungsart} (z.\,B.\ Briefe,
Päckchen, Warensendungen), \emph{Versandart} (z.\,B.\ Einschreiben, Express-Brief,
Dialogpost), \emph{Gewicht} sowie nach \emph{In- und Auslandssendungen}. Diese Sortierung
ist die Grundlage für die spätere Wahl von Format, Frankierung und Dienstleister.

\section{Versand- und Sendungsarten}

Je nach Inhalt, Wert und Eile einer Sendung steht eine Reihe von Versandarten zur
Verfügung. Die Wahl beeinflusst Kosten, Nachweisbarkeit und Sicherheit.

\subsection{Brief und Briefformate}

Briefe werden bei der Deutschen Post nach \emph{Format und Gewicht} eingeteilt. Mit
steigendem Format und Gewicht steigt das Porto.

\begin{center}
\begin{tabular}{p{3cm} p{4.5cm} r}
\toprule
\textbf{Briefformat} & \textbf{Merkmal (Richtwert)} & \textbf{max.\ Gewicht}\\
\midrule
Standardbrief & klein, dünn (z.\,B.\ 1--3 Blatt) & bis 20 g\\
Kompaktbrief & etwas dicker & bis 50 g\\
Großbrief & DIN A4 ungefaltet möglich & bis 500 g\\
Maxibrief & dick / DIN A4 mit Anlagen & bis 1000 g\\
\bottomrule
\end{tabular}
\end{center}

\subsection{Einschreiben}

Möchte man die Zustellung wichtiger Briefe (z.\,B.\ Kündigungen, fristgebundene
Schriftstücke) dokumentieren, bietet sich die Versandart \textbf{Einschreiben} an. Die
Deutsche Post bietet mehrere Varianten:

\begin{itemize}
  \item \textbf{Einschreiben Einwurf} --- der Postzusteller dokumentiert, dass der Brief in den Briefkasten oder das Postfach des Empfängers eingeworfen wurde.
  \item \textbf{Einschreiben (Übergabe)} --- der Zusteller übergibt den Brief gegen Unterschrift an den Empfänger oder einen Empfangsberechtigten (z.\,B.\ Ehepartner); Sendungsverfolgung möglich.
  \item \textbf{Einschreiben Eigenhändig} --- Übergabe nur gegen Unterschrift und ausschließlich an den Empfänger oder einen schriftlich Bevollmächtigten.
  \item \textbf{Einschreiben Rückschein} --- persönliche Übergabe gegen Unterschrift; eine Empfangsbestätigung mit Datum und Unterschrift wird dem Absender zugeschickt.
\end{itemize}

Bei Verlust oder Beschädigung haftet die Deutsche Post bis zur Höhe des unmittelbaren
Schadens, jedoch maximal 25,00~€\ bei Einschreiben bzw.\ 20,00~€\ bei Einschreiben
Einwurf.

\subsection{Einschreiben Wert, Warensendung und Nachnahme}

\textbf{Einschreiben Wert} bietet sich an, wenn Bargeld (bis 100,00~€) oder kleinere
Gegenstände wie Smartphones oder kleine Geschenke (bis 500,00~€) per Brief verschickt
werden sollen. Der Versender erhält einen Einlieferungsnachweis, kann die Sendung online
verfolgen, und die Zustellung wird dokumentiert. Im Haftungsfall erstattet die Deutsche
Post den nachgewiesenen Wert bis max.\ 500,00~€\ bei Sachwerten und max.\ 100,00~€\
bei Bargeld.

Mit der \textbf{Warensendung} lassen sich kleinere Gegenstände und Werbesendungen
(z.\,B.\ Kataloge, Ersatzteile, Zubehör) kostengünstig verschicken. Die Sendungen sind
jedoch \emph{nicht versichert}, und eine Sendungsverfolgung ist nicht möglich.

Bei der \textbf{Nachnahme} wird die Sendung erst gegen Bezahlung des Betrags an den
Empfänger ausgehändigt; der eingezogene Betrag wird anschließend an den Absender
überwiesen --- sinnvoll, wenn vor der Lieferung gezahlt werden soll.

\subsection{Dialogpost}

Die \textbf{Dialogpost} eignet sich für den Versand von Werbebriefen, Kundenmagazinen
oder Einladungen, wenn diese an \emph{mindestens 200 Empfänger derselben Leitregion}
(Übereinstimmung der ersten beiden Ziffern der Postleitzahl), 4\,000 Empfänger bundesweit
oder 500 Empfänger bundesweit (nur werbliche Inhalte) verschickt werden sollen. Sie ist
eine kostengünstige Versandart für Massenwerbung.

\section{Falzen, Kuvertieren und Verschließen}

Das \emph{Falten} (Falzen) der Schriftstücke erfolgt manuell oder mithilfe einer
Falzmaschine. Ziel ist es, für einen Brief das passende, aber auch das kostengünstigste
Format für den Versand zu erhalten. Dafür gibt es verschiedene Falzarten:

\begin{center}
\begin{tabular}{p{4cm} p{4cm} p{4.5cm}}
\toprule
\textbf{Falzart} & \textbf{aus Briefbogen} & \textbf{in Briefhülle}\\
\midrule
Bruch-/Einfachfalz & A4 (Einfachfalz auf A5) & C5 mit Sichtfenster\\
Bruch-/Einfachfalz & A5 (Einfachfalz auf A6) & C6 mit Sichtfenster\\
Kreuzfalz & A4 (Kreuzfalz auf A6) & C6\\
Zickzack-/Leporellofalz & A4 (auf 1/3 A4) & DIN lang mit Sichtfenster\\
Wickelfalz & A4 (auf 1/3 A4) & DIN lang\\
\bottomrule
\end{tabular}
\end{center}

\medskip
Beim \textbf{Kuvertieren} werden die Schriftstücke in die für den Versand geeigneten
Umschläge gesteckt. Dies kann manuell oder maschinell mit Kuvertiermaschinen erfolgen.
Anschließend werden die Briefe \emph{verschlossen}.

\begin{merksatz}
Die richtige Falzart spart bares Geld: Das kleinstmögliche, zulässige Format senkt das
Porto. Ein A4-Brief im Zickzack- oder Wickelfalz passt in eine DIN-lang-Hülle und wird
oft als günstiger Standardbrief versendet.
\end{merksatz}

\section{Wiegen und Frankieren}

Bevor das Porto ermittelt werden kann, muss die Post \emph{gewogen} werden --- mittels
mechanischer oder elektronischer Briefwaage. \emph{Frankieren} bedeutet das Freimachen,
also das Bezahlen des Portos. Dafür stehen verschiedene Frankierarten zur Verfügung:

\begin{itemize}
  \item \textbf{Briefmarken (Postwertzeichen)} --- für einzelne Briefe geeignet.
  \item \textbf{Plusbrief} --- Umschlag mit aufgedruckter Briefmarke.
  \item \textbf{DV-Verfahren} --- Frankieren mittels EDV; muss von der Deutschen Post genehmigt werden und eignet sich bei Massenpost.
  \item \textbf{Frankiermaschine} --- bedruckt den Briefumschlag mit dem Postentgelt, dem Datum und evtl.\ einem Werbeaufdruck. Bei Paketen, Päckchen und sperrigen Sendungen erfolgt der Ausdruck auf einem selbstklebenden Frankierstreifen.
  \item \textbf{Online-Frankierung} --- das Porto wird am PC bezahlt und als Marke/Code ausgedruckt (moderne Praxis).
\end{itemize}

\subsection{Abrechnung des Postentgelts}

Die Abrechnung des Postentgelts bei der Frankiermaschine kann auf zwei Wegen erfolgen:

\begin{itemize}
  \item \textbf{Wertvorgabesystem} --- ein bestimmter Betrag wird im Voraus bezahlt, in der Frankiermaschine eingestellt und durch das Stempeln der Sendungen verbraucht.
  \item \textbf{Fernwertvorgabesystem} --- der Entgelt-Betrag wird per Telefon oder PC beantragt und übermittelt und dann mithilfe eines Codes in die Frankiermaschine eingegeben.
\end{itemize}

\section{Aufgeben der Post und Versanddienstleister}

Ist die Post versandfertig, wird sie \emph{aufgegeben}. Dafür gibt es mehrere Wege:

\begin{itemize}
  \item \textbf{Postfiliale} --- Einlieferung am Schalter.
  \item \textbf{Bring- und Abholservice} --- mit dem Dienst \emph{HIN + WEG} bringt die Deutsche Post morgens die Geschäftspost zu einem vereinbarten Termin ins Unternehmen und holt nachmittags die ausgehenden Sendungen ab.
  \item \textbf{KEP-Dienste} --- Kurier-, Express- und Paketdienste.
\end{itemize}

\subsection{Kurier-, Express- und Paketdienste (KEP)}

\begin{itemize}
  \item \textbf{Kurierdienst} --- befördert die Sendungen persönlich und direkt zum Empfänger; schnell (z.\,B.\ Fahrradkuriere in Großstädten), geeignet für vertrauliche, besonders schutzbedürftige Dokumente oder kleinere Sendungen.
  \item \textbf{Expressdienst} --- übernimmt die Beförderung und \emph{garantiert}, dass die Sendung innerhalb einer bestimmten Frist beim Empfänger eintrifft.
  \item \textbf{Paketdienst} --- übernimmt \emph{keine Garantie} der Laufzeit; die meisten Paketdienste beschränken das Höchstgewicht auf 30~kg.
\end{itemize}

Neben der \emph{Deutschen Post/DHL} bieten zahlreiche \emph{private Dienstleister}
(z.\,B.\ DPD, GLS, UPS, Hermes) Brief- und Paketdienste an --- oft mit eigener
Sendungsverfolgung und Versandsoftware.

\section{Poststraße}

Werden die Maschinen, die für die einzelnen Arbeitsschritte zur Verarbeitung der
ausgehenden Post benötigt werden, hintereinandergeschaltet und in einem Raum zentral
zusammengefasst, so spricht man von einer \textbf{Poststraße} (Kuvertier- bzw.\
Frankierstraße). Sie ermöglicht das \emph{Falzen, Kuvertieren, Verschließen, Trennen nach
Portoklassen und Frankieren} in einem automatischen Maschinengang --- ideal für
Unternehmen mit hohem Postaufkommen.

\subsection{Moderne Praxis: Versandsoftware und Hybridpost}

Heute werden viele dieser Schritte digital ausgelöst und gesteuert:

\begin{itemize}
  \item \textbf{Versandsoftware} ermittelt automatisch die günstigste Versandart und das Porto, erstellt Labels und Frankierungen, verwaltet Adressen und übergibt die Daten elektronisch an den Dienstleister.
  \item \textbf{Hybridpost} --- das Unternehmen übergibt die Briefe nur noch \emph{digital}; ein Dienstleister druckt, kuvertiert, frankiert und stellt sie physisch zu. Der Wechsel von digital zu Papier (Medienbruch) erfolgt erst beim Dienstleister.
\end{itemize}

\section{Elektronische Post}

In vielen Unternehmen erfolgt die Korrespondenz auf elektronischem Wege mittels
\textbf{E-Mails}, die den postalischen Geschäftsbrief ersetzen. Aus diesem Grund müssen
E-Mails bestimmte \emph{Pflichtangaben} enthalten. Dies ist im \emph{EHUG} (Gesetz über
das elektronische Handelsregister und Genossenschaftsregister sowie das
Unternehmensregister) geregelt. Bereits 2003 wurde in der EU-Publizitätsrichtlinie
festgelegt, dass die Pflichtangaben unabhängig davon gelten, ob die geschäftliche
Korrespondenz in Papierform oder in anderer Form erfolgt.

Zu den Geschäftsbriefen zählen auch \emph{Bestellungen, Lieferscheine, Rechnungen und
Quittungen} --- auch in Form von E-Mail oder Telefax. Welche Pflichtangaben nötig sind,
hängt von der \emph{Rechtsform} des Unternehmens ab.

\begin{center}
\begin{tabular}{p{3cm} p{9cm}}
\toprule
\textbf{Rechtsform} & \textbf{Pflichtangaben (Auswahl)}\\
\midrule
Einzelunternehmen & Firma, Bezeichnung nach \S~19 HGB (z.\,B.\ e.\,K.), Ort der Niederlassung, Registergericht, Handelsregisternummer\\
OHG, KG & Firma, Rechtsform, Hauptsitz, Registergericht, Handelsregisternummer\\
GmbH & Firma, Rechtsform, Sitz, Registergericht, Handelsregisternummer, alle Geschäftsführer, ggf.\ Aufsichtsratsvorsitzender\\
AG & Firma, Rechtsform, Hauptsitz, Registergericht, Handelsregisternummer, alle Vorstandsmitglieder, Aufsichtsratsvorsitzender\\
\bottomrule
\end{tabular}
\end{center}

\medskip
Die genaue Adresse, Telefon- oder Faxnummern, ein Postfach oder die Bankverbindungen
zählen \emph{nicht} zu den Pflichtangaben. Soll der E-Mail-Vordruck zugleich als
\emph{Rechnung} dienen, sind nach \S~14 Umsatzsteuergesetz (UStG) zusätzliche Angaben
nötig (z.\,B.\ Umsatzsteuer-Identifikationsnummer).

\subsection{De-Mail, EGVP und elektronische Signatur}

Für \emph{rechtssichere} elektronische Kommunikation gibt es besondere Dienste, die das
klassische Einschreiben im digitalen Verkehr ablösen:

\begin{itemize}
  \item \textbf{De-Mail} --- ein verschlüsselter, nachweisbarer E-Mail-Dienst für verbindliche elektronische Kommunikation.
  \item \textbf{EGVP} (Elektronisches Gerichts- und Verwaltungspostfach) --- rechtsverbindlicher elektronischer Schriftverkehr mit Gerichten und Behörden.
  \item \textbf{Elektronische (qualifizierte) Signatur} --- ersetzt die eigenhändige Unterschrift, weist die Identität des Absenders nach und sichert, dass das Dokument nach dem Signieren nicht verändert wurde (Integrität).
\end{itemize}

\section{Zusammenfassung}

\begin{enumerate}
  \item Der Postausgang läuft in \textbf{sieben Schritten} ab: Kontrollieren, Sortieren, Falzen, Kuvertieren/Verschließen, Wiegen, Frankieren, Aufgeben.
  \item \textbf{Versandarten} werden nach Inhalt, Wert und Eile gewählt: Brief (Standard/Kompakt/Groß/Maxi), Einschreiben (Einwurf/Übergabe/Eigenhändig/Rückschein), Einschreiben Wert, Warensendung, Dialogpost, Nachnahme.
  \item \textbf{Falzarten} (Einfach-, Kreuz-, Zickzack-, Wickelfalz) bringen den Brief ins kostengünstigste Format.
  \item \textbf{Frankierarten}: Briefmarke, Plusbrief, DV-Verfahren, Frankiermaschine, Online-Frankierung; Abrechnung per Wertvorgabe- oder Fernwertvorgabesystem.
  \item Die \textbf{Poststraße} automatisiert mehrere Schritte; \textbf{Versandsoftware} und \textbf{Hybridpost} sind die moderne Weiterentwicklung.
  \item \textbf{Elektronische Post}: E-Mail mit Pflichtangaben nach EHUG; \textbf{De-Mail, EGVP und elektronische Signatur} schaffen Rechtssicherheit.
\end{enumerate}

\newpage

\section*{Glossar}
\addcontentsline{toc}{section}{Glossar}

\glossareintrag{Poststelle}{Zentrale Stelle im Unternehmen, in der die gesammelte Ausgangspost versandfertig gemacht wird.}

\glossareintrag{Kontrollieren}{Prüfung der Ausgangspost auf Unterschrift, Adresse und Vollständigkeit der Anlagen.}

\glossareintrag{Sortieren}{Ordnen der Post nach Sendungsart, Versandart, Gewicht sowie In- und Auslandssendungen.}

\glossareintrag{Fensterbriefumschlag}{Umschlag mit Sichtfenster; die Empfängeradresse wird durch das Fenster sichtbar, ein erneutes Adressieren entfällt.}

\glossareintrag{Einschreiben}{Versandart zur dokumentierten Zustellung wichtiger Briefe; Varianten: Einwurf, Übergabe, Eigenhändig, Rückschein.}

\glossareintrag{Einschreiben Wert}{Versandart für Bargeld (bis 100,00 €) oder Sachwerte (bis 500,00 €) mit Einlieferungsnachweis und Sendungsverfolgung.}

\glossareintrag{Warensendung}{Kostengünstige Versandart für kleinere Gegenstände und Werbesendungen; nicht versichert, keine Sendungsverfolgung.}

\glossareintrag{Dialogpost}{Versandart für Massenwerbung an mind.\ 200 Empfänger derselben Leitregion bzw.\ 4\,000/500 Empfänger bundesweit.}

\glossareintrag{Nachnahme}{Versandart, bei der die Sendung erst gegen Bezahlung des Betrags an den Empfänger ausgehändigt wird.}

\glossareintrag{Falzen}{Falten der Schriftstücke ins passende, kostengünstige Format; Falzarten: Einfach-, Kreuz-, Zickzack-/Leporello-, Wickelfalz.}

\glossareintrag{Kuvertieren}{Einstecken der Schriftstücke in geeignete Umschläge, manuell oder mit Kuvertiermaschinen.}

\glossareintrag{Frankieren}{Freimachen einer Sendung, also Bezahlen des Portos.}

\glossareintrag{Frankiermaschine}{Maschine, die den Umschlag mit Postentgelt, Datum und evtl.\ Werbeaufdruck bedruckt.}

\glossareintrag{Wertvorgabesystem}{Im Voraus bezahlter Betrag wird in der Frankiermaschine eingestellt und durch Stempeln verbraucht.}

\glossareintrag{Fernwertvorgabesystem}{Der Entgelt-Betrag wird per Telefon/PC beantragt und per Code in die Frankiermaschine eingegeben.}

\glossareintrag{KEP-Dienste}{Kurier-, Express- und Paketdienste; befördern Sendungen schnell, fristgarantiert oder als Paket (bis 30 kg).}

\glossareintrag{Poststraße}{Hintereinandergeschaltete Maschinen, die Falzen, Kuvertieren, Verschließen, Trennen und Frankieren automatisch ausführen.}

\glossareintrag{Hybridpost}{Digital übergebene Briefe werden vom Dienstleister gedruckt, kuvertiert, frankiert und physisch zugestellt.}

\glossareintrag{EHUG}{Gesetz über das elektronische Handels-, Genossenschafts- und Unternehmensregister; regelt die Pflichtangaben in E-Mails.}

\glossareintrag{De-Mail}{Verschlüsselter, nachweisbarer E-Mail-Dienst für verbindliche elektronische Kommunikation.}

\glossareintrag{EGVP}{Elektronisches Gerichts- und Verwaltungspostfach für rechtsverbindlichen Schriftverkehr mit Gerichten und Behörden.}

\glossareintrag{Elektronische Signatur}{Ersetzt die eigenhändige Unterschrift; weist die Identität nach und sichert die Unverändertheit des Dokuments.}

\end{document}
